\capitulo{5}{Resultados}

\section{Resumen de resultados.}
En relación con los objetivos planteados al principio de este documento, los resultados obtenidos con el desarrollo de la plataforma propuesta son los siguientes:
\begin{itemize}
    \item \textbf{Desarrollar una aplicación web interactiva con Shiny}:Se implementó una app Shiny con mapas de mortalidad y tasas de incidencia provinciales por melanoma y mapas con variables meteorologicas, también se incluyo historicos de estos datos para futuras líneas entre otras funcionalidades, desplegada en Shiny Server y accesible públicamente.\ref{fig:portada}
     \begin{figure}[H]
        \centering
        \includegraphics[scale=1]{img/portada.png}
        \caption{Pestaña inicial UV-Derma. Elaboración Porpia}
        \label{fig:portada}
    \end{figure}
    \item \textbf{Dotar al usuario de conocimientos sobre el impacto de variables ambientales en salud}:Se incorporó un panel informativo interactivo con descripciones sobre el melanoma, también se efectos mencionan los efectos de los factores de riesgo mencionados en la plataforma. \ref{fig:info1}, \ref{fig:info}
     \begin{figure}[H]
        \centering
        \includegraphics[scale=1]{img/Captura de pantalla 2025-05-28 005537.png}
        \caption{Pestaña Contenido Informativo. Elaboración Porpia}
        \label{fig:info}
    \end{figure}
    \begin{figure}[H]
        \centering
        \includegraphics[scale=1]{img/info1.png}
        \caption{Pestaña Marco Teorico. Elaboración Porpia}
        \label{fig:info1}
    \end{figure}
    \item \textbf{Recomendar al usuario medidas preventivas del melanoma segun los factores de riesgo en una zona geográfica y fototipo de piel determinado}: Se desarrolló un módulo de recomendaciones geolocalizadas que, en función de los valores de radiación UV, temperatura de la ubicación y fototipo del usuario seleccionado, sugiere medidas de protección.\ref{fig:alerta}
     \begin{figure}[H]
        \centering
        \includegraphics[scale=1]{img/recomendacion.png}
        \caption{Pestaña Recomendador. Elaboración Porpia}
        \label{fig:alerta}
    \end{figure}
    \item \textbf{Uso y limpieza de datos abiertos (INE y AEMET)}: Se limpiaron y procesaron los datos meteorológicos y de mortalidad para adaptarlos a la generación de mapas interactivos y su posterior consumo por otras funciones de la aplicación. \ref{fig:tasa_incidencias} \ref{fig:temperaturas}, \ref{fig:Indices_UV}
     \begin{figure}[H]
        \centering
        \includegraphics[width=0.8\textwidth]{img/tasaincidencias.png}
        \caption{Pestaña Mapa defunciones INE. Elaboración Porpia}
        \label{fig:tasa_incidencias}
    \end{figure}
     \begin{figure}[H]
        \centering
        \includegraphics[scale=1]{img/temp.png}
        \caption{Pestaña Mapa Temperaturas. Elaboración Porpia}
        \label{fig:temperaturas}
    \end{figure}
     \begin{figure}[H]
        \centering
        \includegraphics[scale=1]{img/mapauv.png}
        \caption{Pestaña Mapa Índices UV. Elaboración Porpia}
        \label{fig:Indices_UV}
    \end{figure}
    \item \textbf{Facilitar una estimación aproximada del riesgo acumulado de sufrir melanoma}: Se implementó un módulo de cálculo de riesgo acumulado basado en la zona geofráfica y factores personales que devuelve una estimación orientativa del riesgo acumulado de desarrollar melanoma.
    \ref{fig:Riesgo}
     \begin{figure}[H]
        \centering
        \includegraphics[width=0.8\textwidth]{img/riesgo.png}
        \caption{Pestaña Riesgo Acumulado. Elaboración Porpia}
        \label{fig:Riesgo}
    \end{figure}
    \item \textbf{Contribuir al ámbito de la epidemiología ambiental}: Esta plataforma se ha puesto a disposición del público fomentando una educación y prevención del melanoma. También estará disponible todo su desarrollo en un repositorio GitHub si se quiere profundizar en la investigación en las herramientas proporcionadas.

\end{itemize}
 



\clearpage
\section{Discusión}
Viendo el apartado anterior se puede observar que los resultados obtenidos reflejan una correspondecia directa con objetivos propuestos al inicio de este proyecto, ofreciendo tanto un desarrollo técnico y atractivo  como una herramienta de divulgación y prevención accesible para el melanoma cutáneo.

El proceso de limpieza y preparación de los datos procedentes de INE y AEMET garantizó la fiabilidad de las visualizaciones y el correcto funcionamiento de las llamadas a API. Sin embargo, la dependencia de servicios externos introduce vulnerabilidades ante posibles cambios en los \textit{endpoints} \footnote{En desarrollo de software, un \textit{endpoint} es la dirección URL o punto de acceso de una API o servicio web al que la aplicación envía solicitudes para obtener o enviar datos.} lo cual debe tenerse en cuenta en la estabilidad actual del sistema.

El módulo de recomendaciones geolocalizadas ha mostrado ser efectivo para comunicar medidas preventivas adaptadas al índice UV y la temperatura local. Por su parte, el cálculo del índice de riesgo acumulado ofrece una valoración clara y orientativa, facilitando su comprensión por el público general. No obstante, al tratarse de una estimación aproximada diseñada para simplificar la interpretación de riesgos poco perceptibles, no sustituye una evaluación clínica. 

Por otro lado, esta plataforma cumple con los requisitos del GDPR ya que no se trata en ningun momento con datos personales del usuario.

Finalmentee, como contribucion a la epidemiología ambiental se dispone de un repositorio público como recurso reutilizable por desarrolladores e investigadores para mejorar las metodologías de prevencion y asistencia sanitaria.

En resumen,UV-Derma demuestra su valía como herramienta de apoyo a la prevención del melanoma, cumpliendo los objetivos de desarrollo, divulgación y accesibilidad. 


